\documentclass[11pt]{article}
\usepackage[margin=1in]{geometry}
\usepackage{amsmath,amssymb}
\usepackage{booktabs}
\usepackage{enumitem}
\usepackage{hyperref}
\usepackage{xcolor}
\usepackage{listings}

\hypersetup{colorlinks=true, linkcolor=blue, urlcolor=blue}
\emergencystretch=2em

\lstset{
  basicstyle=\ttfamily\small,
  backgroundcolor=\color{gray!10},
  frame=single,
  framerule=0pt,
  breaklines=true,
  language=Python,
  showstringspaces=false
}

\newcommand{\E}{\mathbb{E}}
\newcommand{\Var}{\mathrm{Var}}
\newcommand{\Cov}{\mathrm{Cov}}
\newcommand{\ATE}{\mathrm{ATE}}
\newcommand{\ATT}{\mathrm{ATT}}
\newcommand{\indep}{\perp\!\!\!\perp}
\newcommand{\points}[1]{\hfill \textbf{[#1 points]}}

\title{Assignment 1: Causal Questions, Regression, and Selection on Observables}
\author{Applied Econometrics}
\date{}

\begin{document}
\maketitle

\noindent\textbf{Total: 15 points.} This assignment is based on Lectures 1 and 2. It mixes short theory questions with a small Python exercise. The aim is not to produce a polished empirical paper. The aim is to show that you can define a causal question, write the relevant estimand, and understand what regression is doing.

\paragraph{Submission.} Submit one PDF write-up and one Python script or notebook. Your write-up should include the requested tables, coefficients, and short explanations. Unless a question asks for a derivation, two to five clear sentences are enough.

\paragraph{Starter data.} Run the following command from the course root:
\begin{lstlisting}
python assignments/code/assignment-1-starter.py
\end{lstlisting}
It creates:
\begin{itemize}
  \item \texttt{assignments/data/assignment1\_observational.csv}
  \item \texttt{assignments/data/assignment1\_randomized.csv}
\end{itemize}

\paragraph{Variables.}
\begin{itemize}
  \item Outcome: \texttt{final\_score}.
  \item Treatment: \texttt{scholarship}.
  \item Pre-treatment controls: \texttt{female}, \texttt{baseline\_score}, \texttt{parent\_edu}, and \texttt{age}.
  \item Post-treatment variable: \texttt{study\_hours\_after}.
\end{itemize}

\section*{Part A: Causal Setup and Selection Bias}

\begin{enumerate}[label=A\arabic*.]
  \item Consider the question: ``What is the effect of receiving a study scholarship on a student's final exam score?'' Define \(Y_{1i}\), \(Y_{0i}\), \(D_i\), the observed outcome \(Y_i\), and the average treatment effect \(\ATE\). State clearly what the counterfactual is. \points{1.5}

  \item Show the selection-bias decomposition:
  \[
    \E[Y_i\mid D_i=1]-\E[Y_i\mid D_i=0]
    =
    \E[Y_{1i}-Y_{0i}\mid D_i=1]
    +
    \left\{\E[Y_{0i}\mid D_i=1]-\E[Y_{0i}\mid D_i=0]\right\}.
  \]
  Explain in words what each term means. \points{2}

  \item Suppose more motivated students are more likely to apply for and receive the scholarship. In the regression
  \[
    Y_i=\alpha+\rho D_i+u_i,
  \]
  explain why \(\rho\) may be biased. Use the omitted-variable bias formula to predict the likely sign of the bias if motivation raises scores and raises scholarship take-up. \points{1.5}
\end{enumerate}

\section*{Part B: Coding Exercise}

Use the observational and randomized datasets created by the starter script. Report coefficients to three decimal places.

\begin{enumerate}[label=B\arabic*.]
  \item In \texttt{assignment1\_observational.csv}, estimate the naive difference in means:
  \[
    \bar{Y}_{D=1}-\bar{Y}_{D=0}.
  \]
  Then estimate the same object by OLS using \(Y_i=\alpha+\rho D_i+u_i\). Verify that the two estimates are identical up to rounding. \points{1}

  \item Make a small balance table comparing treated and untreated students on \texttt{baseline\_score}, \texttt{parent\_edu}, \texttt{age}, and \texttt{female}. Which variables look most imbalanced? Explain why this matters for the causal interpretation of the naive comparison. \points{1}

  \item Estimate
  \[
    Y_i=\alpha+\rho D_i+X_i'\gamma+u_i,
  \]
  where \(X_i\) includes the four pre-treatment controls. Compare \(\hat{\rho}\) to the naive estimate. Explain the movement using the selection-bias or omitted-variable-bias logic from Part A. \points{2}

  \item Add \texttt{study\_hours\_after} to the controlled regression. Does the scholarship coefficient change? Should this be your preferred estimate of the total scholarship effect? Explain why this variable is a bad-control candidate. \points{1}

  \item Repeat B1 and B3 using \texttt{assignment1\_randomized.csv}. Compare the naive and controlled estimates. What should random assignment do to the need for controls? \points{1}
\end{enumerate}

\section*{Part C: Short Applied Interpretation}

\begin{enumerate}[label=C\arabic*.]
  \item State the conditional independence assumption needed to interpret the controlled observational regression causally. Write it in symbols and then in plain English. \points{1.5}

  \item State the overlap/common-support condition in words. Give one simple diagnostic from your data that helps assess overlap or comparability. \points{1}

  \item Write one short paragraph giving your preferred estimate from this assignment and the main remaining threat to a causal interpretation. Your paragraph should mention whether your preferred evidence comes from the observational or randomized dataset. \points{1.5}
\end{enumerate}

\section*{Point Summary}
\begin{center}
\begin{tabular}{lcc}
\toprule
Part & Topic & Points \\
\midrule
A & Potential outcomes, selection bias, OVB & 5 \\
B & Python regression and balance checks & 6 \\
C & Causal interpretation & 4 \\
\midrule
Total & & 15 \\
\bottomrule
\end{tabular}
\end{center}

\end{document}
